\documentclass[conference]{IEEEtran}
\usepackage{amsmath}
\usepackage{amssymb}
\usepackage{graphicx}
\usepackage{booktabs}
\usepackage{xcolor}
\usepackage{hyperref}
\usepackage{cite}
\usepackage{float}

\hypersetup{colorlinks=true, linkcolor=blue, urlcolor=blue, citecolor=blue}

\title{Proprioception as a Segmentation Prompt: Force-Guided Object Grounding for Learning from Demonstration}

\author{
\IEEEauthorblockN{Anurag Shetye}
\IEEEauthorblockA{University of Twente / NAKAMA Robotics Lab (remote collaboration)}
}

\begin{document}
\maketitle

\begin{abstract}
Learning-from-demonstration (LfD) pipelines that reason about kinematic
constraints or physical interaction depend on knowing which object in the
scene is relevant at each phase of a demonstration. This is typically
solved either by hand-labelling, by gripper-event heuristics that only
identify the grasped object, or by label-free visual object discovery. We
show empirically, across multiple real demonstrations, that label-free
visual discovery (DINOv2 self-attention combined with monocular or real
sensor depth) fails to isolate task-relevant objects in typical in-lab
manipulation scenes -- including when real, not estimated, depth is
available -- indicating that vision alone is insufficient. We propose
using the robot's own proprioception as the grounding signal instead:
gripper-state transitions localise the grasped object in time, and the
wrist force/torque wrench, via the intrinsic contact ray of Bicchi
\cite{bicchi1990}, geometrically constrains the location of the
contact-receiving object once camera calibration is available. Both cues
seed a frozen promptable video segmentation model (SAM~2
\cite{ravi2024sam2}) and propagate a role-tagged object mask across the
full demonstration. We report a working system validated on four
real-robot recordings spanning three distinct manipulation tasks and
three different sensor-availability profiles (full sensing, force-only,
gripper-only), including a first case of tracking three simultaneous
task-relevant objects through one demonstration. The geometric wrench-ray
grounding experiment -- the paper's core evaluation -- is fully
implemented and awaiting camera calibration data from the host lab; we
report the evaluation protocol and preliminary (dry-run, uncalibrated)
evidence that the signal is directionally correct.
\end{abstract}

\section{Introduction}

Learning from demonstration increasingly decomposes a demonstration into
phases governed by kinematic constraints \cite{overbeek2025frontiers} or
physical contact \cite{overbeek2025adrr}. Both layers presuppose an
answer to a prior question: \emph{which object, in the camera frame, is
the one the constraint or contact concerns?} Most existing approaches
either sidestep this (constraints defined on end-effector pose alone), or
identify only the grasped object from gripper state, leaving any
contact-receiving object unidentified.

A natural alternative is label-free visual object discovery: use a
frozen foundation model to propose salient objects without any
task-specific training. We test this directly, using DINOv2
\cite{oquab2023dinov2} self-attention combined with monocular depth
(Depth-Anything-V2 \cite{yang2024depthanythingv2}) and, in a second test,
combined with \emph{real} sensor depth from an RGB-D camera. In both
cases the resulting saliency fails to isolate the task-relevant object in
our recordings (Section~\ref{sec:experiments}), including on non-specular
objects where depth is measured cleanly. This motivates a proprioceptive
alternative.

\textbf{Contributions.}
\begin{itemize}
  \item An event-detection scheme that localises grasp, release, and
        contact moments from gripper and wrist force/torque signals,
        with graceful degradation when either sensor is unavailable
        (verified on real recordings missing each sensor in turn).
  \item A geometric grounding method for the contact-receiving object,
        using the intrinsic contact ray \cite{bicchi1990} computed from
        the wrist wrench and projected into the camera frame via
        calibrated extrinsics -- implemented and validated in dry-run
        (uncalibrated) mode, pending calibration data to report
        quantitative accuracy.
  \item A generalisation of the resulting role-tagged tracking pipeline
        from a fixed grasped/contact-receiver object pair to an
        arbitrary number of simultaneously tracked objects, demonstrated
        with three independently seeded and propagated objects in one
        recording.
  \item An empirical case, using real (not estimated) depth, that
        label-free visual object discovery is not a substitute for
        proprioceptive grounding in this setting.
\end{itemize}

\section{Related Work}

\textbf{Gripper-event affordance learning.} Prior self-supervised
approaches use gripper open/close transitions to bootstrap object
affordance labels, but do not address which object receives contact once
grasped.

\textbf{Label-free / promptable segmentation.} Segment Anything
\cite{kirillov2023sam} and its video extension SAM~2
\cite{ravi2024sam2} provide high-quality, class-agnostic segmentation
given a point or box prompt, but supply no mechanism for choosing
\emph{where} to prompt. Depth-and-attention object discovery methods
(DADO-style: DINOv2 attention re-weighted by monocular depth) attempt
label-free prompt selection; we show this fails on our scenes even with
real depth substituted for the monocular estimate.

\textbf{Contact sensing.} Bicchi's intrinsic contact sensing
\cite{bicchi1990} recovers the point and direction of a single contact
from a wrist force/torque measurement alone, without vision. To our
knowledge this geometric signal has not previously been used as a prompt
for visual segmentation.

\textbf{Constraint and interaction identification.} The host lab's
existing LfD layers identify kinematic constraints from pose
\cite{overbeek2025frontiers} and physical interactions from contact
\cite{overbeek2025adrr}; this work is intended as the visual grounding
layer beneath both.

\section{Method}

\subsection{Event detection}

Given a merged, pose-synchronised recording, three event types are
detected: \emph{grasp} and \emph{release} from gripper-width threshold
crossings (midpoint of the observed open/closed range), and
\emph{contact} (``press'') from the baseline-subtracted wrist
force-magnitude peak. When only one of the two sensors is present, the
detector degrades gracefully to that sensor's events alone (verified on
real recordings lacking each sensor in turn, Section~\ref{sec:experiments}).
For recordings with more than one grasp/release/contact cycle, peaks are
grouped into activity clusters by time-gap rather than reporting only the
single global event.

\subsection{Geometric grounding of the contact-receiving object}

At a contact event, the wrist wrench $(\mathbf{f}, \boldsymbol{\tau})$
measured in the sensor frame determines, under a single-contact
assumption, the contact ray's closest point to the sensor origin and its
direction:
\begin{equation}
  \mathbf{r}_0 = \frac{\mathbf{f} \times \boldsymbol{\tau}}{\lVert \mathbf{f} \rVert^2}, \qquad
  \hat{\mathbf{f}} = \frac{\mathbf{f}}{\lVert \mathbf{f} \rVert}
  \label{eq:bicchi}
\end{equation}
following Bicchi \cite{bicchi1990}. Transforming $\mathbf{r}_0$ and
$\hat{\mathbf{f}}$ through the sensor-to-base and base-to-camera
transforms and projecting through the calibrated camera intrinsics
yields a 2D ray in the image, which is used as a point (or line) prompt
for SAM~2 on the contact-receiving object -- geometrically, without any
learned object proposal step.

\subsection{Grasped-object grounding}

The grasped object is seeded either by a vision-only heuristic (SAM
automatic mask generation, scored by position/size/border-distance
priors) or, once available, by projecting the end-effector position
directly into the image via the same calibration. The vision-only
heuristic is shown in Section~\ref{sec:experiments} to be unreliable
across scenes with different object scales, motivating the geometric
alternative.

\subsection{Propagation and multi-object generalisation}

Each seeded object is propagated bidirectionally through the full
recording with SAM~2's video predictor, yielding a per-frame mask. The
sidecar representation generalises to an arbitrary number of
independently seeded objects rather than a fixed grasped/contact-receiver
pair, each with its own role label.

\section{System and Data}

Four real-robot recordings on a Franka Research 3 arm with a Franka Hand
gripper, Bota wrist force/torque sensor, and ZED RGB(-D) camera are used.
Table~\ref{tab:recordings} summarises sensor availability and scene
content; the three recordings differ in which sensors were actively
publishing, exercising the graceful-degradation behaviour of the event
detector under real (not synthetic) conditions.

\begin{table}[t]
\centering
\caption{Recordings used in this work.}
\label{tab:recordings}
\begin{tabular}{lccl}
\toprule
\textbf{Recording} & \textbf{Gripper} & \textbf{F/T} & \textbf{Task} \\
\midrule
1 & \checkmark & \checkmark & pick / press / drop (2 objects) \\
2 & --         & \checkmark & multi-object contact (4 objects) \\
3 & \checkmark & --         & fastener insertion \\
4 & \checkmark & --         & fastener insertion \\
\bottomrule
\end{tabular}
\end{table}

Recordings 2--4 include registered depth in addition to RGB, decoded from
the raw sensor transport where the standard export tooling could not
process it.

\section{Experiments}
\label{sec:experiments}

\subsection{Label-free visual discovery fails, even with real depth}

DINOv2 attention combined with monocular depth (Depth-Anything-V2) was
evaluated on Recording~1's event frames and fails to isolate either
task-relevant object -- attention is low and spatially uniform on these
in-lab scenes. To rule out monocular depth estimation error as the cause,
the same recipe was re-run on Recording~2 substituting real sensor depth,
across five contact events. Depth-sensor validity itself varies sharply
with surface material (16\% valid pixels on a specular object vs.\
58--72\% on matte objects in the same scene), and the resulting object
proposal correctly localises the target object in 2 of 5 events, hitting
background clutter in the remaining 3. Real depth is measurably better
than a monocular estimate but does not make label-free discovery
reliable.

\begin{figure}[t]
\centering
\includegraphics[width=\columnwidth]{figures/dado_real_depth_multi.png}
\caption{Label-free discovery (DINOv2 attention $\times$ real depth) on
five contact events. Object proposal (rightmost column, green) correctly
localises the object in the first two rows; the remaining three land on
background surfaces.}
\label{fig:dado}
\end{figure}

\subsection{Vision-only seed placement is unreliable across scenes}

The vision-only heuristic used to seed the grasped/contact-receiver roles
scores SAM's candidate masks by position, size, and border distance, with
a fixed area cutoff tuned on Recording~1's object scale. On Recording~2,
the correct contact-receiving object exceeds this cutoff and is rejected
outright, causing the seed to fall on background clutter and propagate an
unstable mask (17--78\% of frame area) for the full recording. Two
automated generalisations of the heuristic were attempted -- relative
candidate-size ranking, and a real-depth near-camera prior -- and both
failed: the first mis-selects a reflective background surface on
Recording~1 that it correctly avoided under the original heuristic; the
second is blind precisely on the specular surfaces where the original
failure occurs, since sensor depth dropout is worst there. This negative
result motivates the geometric grounding method of
Section~\ref{sec:experiments}C over further heuristic tuning.

\subsection{Wrench-ray geometric grounding (protocol and preliminary evidence)}

This is the paper's core planned evaluation and is currently blocked on
camera intrinsics and extrinsics from the host lab; the method itself
(Eq.~\ref{eq:bicchi} through image projection) is fully implemented and
runs in a dry-run mode that reports base-frame geometry without
projecting, so that no placeholder calibration is ever silently used.
Preliminary (uncalibrated) evidence: at the contact event of
Recording~1, wrist wrench magnitude was 26.4\,N directed within
$5^\circ$ of vertical in the base frame, consistent with a top-down
press onto the visible contact-receiving object; the Bicchi contact
point $\mathbf{r}_0$ lies within 1\,cm of the wrist origin along the
force axis, a physically plausible contact offset for the demonstrated
task.

\textbf{Planned evaluation.} Once calibration is available, for each
contact event across all recordings with real depth (Recordings 2--4):
project $\mathbf{r}_0$ and $\hat{\mathbf{f}}$ into the image, intersect
the ray with the scene using the recorded depth map, and measure the 3D
distance between the intersection point and the nearest point on the
true contact-receiving object's surface. We expect this to report a
single, low (sub-centimetre to few-centimetre) error figure per event as
the headline quantitative result.

\subsection{Generalisation to multiple simultaneous objects}

The fixed-pair role model was generalised to an arbitrary object count
and demonstrated on Recording~2 with three independently seeded and
bidirectionally propagated objects (Fig.~\ref{fig:multiobj}), each
tracked across the full recording (1013/1013 frames) with no
cross-object mask contamination and no unbounded mask growth.

\begin{figure}[t]
\centering
\includegraphics[width=\columnwidth]{figures/identify_depth_multi/overlays/f0786_1782835535683396106.png}
\caption{Three independently-seeded, independently-propagated objects
composed into one sidecar frame, generalising past a fixed
grasped/contact-receiver pair.}
\label{fig:multiobj}
\end{figure}

\subsection{Sensor-degradation robustness}

Event detection was verified to degrade gracefully and correctly across
all three sensor-availability profiles present in the four recordings:
full sensing (Recording~1), force-only (Recording~2, no gripper events
detected or attempted), and gripper-only (Recordings 3--4, no contact
event detected or attempted). Each fallback path was regression-tested
against a recording with full sensing to confirm zero behavioural change
when both sensors are present.

\section{Limitations}

All recordings come from a single lab, a single robot arm, and a small
number of camera poses; generalisation across hardware setups is
untested. The sample size (four recordings, three tasks) is small for a
generalisation claim and is treated here as a system validation rather
than a statistically powered study. The paper's central quantitative
claim -- wrench-ray grounding accuracy -- is not yet reported, pending
calibration data; this is disclosed rather than approximated with
placeholder values. Label-free discovery was tested with one
attention-based method and one depth-conditioning scheme; other label-free
approaches may perform differently.

\section{Conclusion and Future Work}

Proprioceptive signals -- gripper state and wrist wrench -- provide a
grounding signal for task-relevant object identification that label-free
visual discovery does not reliably replace, including when real depth is
available. A geometric method for grounding the contact-receiving object
from wrist force/torque alone is implemented and ready for evaluation
once camera calibration is available; this is the immediate next step.
Beyond that, an integration of the resulting role-tagged object tracks
into the host lab's existing kinematic-constraint and physical-interaction
identification layers \cite{overbeek2025frontiers,overbeek2025adrr} is
the intended path to a complete perception-to-LfD pipeline.

\begin{thebibliography}{9}

\bibitem{bicchi1990}
A.~Bicchi, ``Intrinsic contact sensing for soft fingers,'' in
\emph{Proc. IEEE Int. Conf. Robotics and Automation (ICRA)}, 1990.

\bibitem{kirillov2023sam}
A.~Kirillov \emph{et al.}, ``Segment Anything,'' in \emph{Proc. IEEE/CVF
Int. Conf. Computer Vision (ICCV)}, 2023.

\bibitem{ravi2024sam2}
N.~Ravi \emph{et al.}, ``SAM 2: Segment Anything in Images and Videos,''
\emph{arXiv preprint arXiv:2408.00714}, 2024.

\bibitem{oquab2023dinov2}
M.~Oquab \emph{et al.}, ``DINOv2: Learning Robust Visual Features without
Supervision,'' \emph{arXiv preprint arXiv:2304.07193}, 2023.

\bibitem{yang2024depthanythingv2}
L.~Yang \emph{et al.}, ``Depth Anything V2,'' \emph{arXiv preprint
arXiv:2406.09414}, 2024.

\bibitem{overbeek2025frontiers}
W.~Overbeek \emph{et al.}, ``Kinematic constraint identification for
learning from demonstration,'' \emph{Frontiers in Robotics and AI}, 2025.

\bibitem{overbeek2025adrr}
W.~Overbeek \emph{et al.}, ``Physical interaction identification from
contact for learning from demonstration,'' \emph{Advanced Robotics
Research}, 2025.

\end{thebibliography}

\end{document}
